% !TEX root = ../main.tex

\chapter{结论与展望}

\section{工作总结}

本文围绕听障交流与训练辅助场景，
设计并实现了 HearBridge 多端协同系统。

系统以听障用户的手语学习、口语训练、手势识别和交流辅助为主要应用目标，
构建了 HarmonyOS 移动端、Spring Boot 后端、Vue 后台管理端和 Python 手势识别服务协同工作的整体架构。

在系统需求分析阶段，
本文从普通用户和管理员两个角色出发，
梳理了系统在课程学习、训练交互、数字人动作展示、手势实时识别、句子视频识别、资源管理、样本管理和模型版本管理等方面的功能需求，
并分析了系统在性能、安全性、易用性和可维护性等方面的非功能性需求。

在系统设计阶段，
本文采用多端分离的方式组织系统结构。

移动端主要负责用户交互、相机采集、训练流程展示和识别结果呈现；
后端主要负责业务数据管理、用户认证、文件资源管理、训练记录管理、模型版本管理以及与 Python 识别服务的交互；
后台管理端主要负责手势分类、手势资源、样本数据、训练任务和模型版本的维护；
Python 识别服务主要负责 MediaPipe 关键点检测、特征提取、模型训练、模型加载、实时识别和句子视频识别。

这种架构使系统各端职责相对清晰，
便于后续独立维护和扩展。

在系统实现阶段，
本文完成了移动端、后端、管理端和识别服务的主要功能开发。

移动端实现了登录态恢复、首页浏览、训练入口、手势资源列表、数字人动作展示、全屏相机实时识别、raw 样本采集和句子视频识别结果展示等功能。

后台管理端实现了管理员登录、路由守卫、手势分类管理、手势资源管理、样本管理、训练管理和模型版本管理等功能。

后端实现了用户认证、管理员认证、资源接口、文件上传、样本同步、模型训练触发、模型版本发布和语义修正调用等功能。

Python 识别服务实现了实时单词识别 WebSocket、raw dataset 采集、特征转换、模型训练、模型运行时管理、模型重载和句子视频上传识别等功能。

在手势识别服务实现与实验分析方面，
本文首先针对移动端实时识别场景，
构建了基于 MediaPipe 关键点和时序特征的固定词表识别流程。

系统能够从移动端采集图像帧，
经 Python 服务完成关键点检测、特征提取和模型推理，
再将识别标签、置信度和模型版本信息返回移动端展示。

同时，
系统支持通过移动端采集 raw dataset 样本，
并在后台管理端完成样本同步、特征转换、模型训练、模型版本发布和运行时重载，
形成了从数据采集到模型发布再到实时识别的完整闭环。

针对句子视频识别场景，
本文基于 WLASL 小词表数据构建了词级分类模型，
并采用滑动窗口方式将短句视频转换为多个固定长度的候选窗口。

系统通过逐窗口分类、词段合并、过滤和 NMS 抑制生成原始 Gloss 序列，
再由后端基于词段 Top-K 候选生成有限候选序列，
并调用 DeepSeek 进行受约束的候选选择和自然化表达。

实验结果表明，
候选序列生成能够显著提高理论可恢复空间，
DeepSeek 受约束语义修正在当前小词表条件下能够带来稳定但有限的提升。

同时，
实验也说明当前句子视频识别仍受到词表规模、训练资源、动作边界和候选排序能力的限制。

在系统测试方面，
本文对 HearBridge 系统进行了白盒测试、黑盒功能测试和非功能性测试。

白盒测试覆盖了后端 Controller 层、Service 层和 Interceptor 层的核心逻辑，
自动化测试用例全部通过，
能够支撑对后端关键业务逻辑和鉴权逻辑的正确性验证。

黑盒功能测试从移动端、后台管理端和跨端联动流程出发，
验证了用户训练、数字人展示、实时识别、句子视频识别、资源维护、样本管理、模型训练和模型发布等功能流程。

非功能性测试进一步验证了系统在性能、稳定性、安全性、易用性和可维护性方面的表现。

JMeter 性能测试结果表明，
在 20 并发和 50 并发条件下，
后端核心接口错误率均为 0.00\%，
平均响应时间和高百分位响应时间保持在较低水平，
能够满足毕业设计演示环境下的基本并发访问需求。

总体来看，
本文完成了 HearBridge 系统从需求分析、系统设计、功能实现、识别实验到系统测试的完整工作。

系统已经能够支持听障用户进行手语训练、数字人动作学习、实时手势识别和句子视频识别，
也能够支持管理员完成手势资源、样本数据、训练任务和模型版本的后台维护。

虽然当前系统距离开放域连续手语识别和真实场景大规模应用仍有差距，
但已经实现了一个功能相对完整、链路能够闭环、具备继续扩展基础的听障交流与训练辅助系统。

\section{未来展望}

HearBridge 系统目前已经完成毕业设计阶段的主要功能目标，
但受数据资源、模型能力、工程周期和实验条件限制，
系统仍有进一步优化空间。

后续工作可以从以下几个方面展开。

首先，
需要继续扩充高质量手语数据资源。

当前系统中的词级识别和句子视频识别仍主要依赖有限词表和小规模样本。

由于高质量手语视频资源获取难度较高，
部分动作类别的训练样本数量有限，
模型在相似手型、相似姿态和相似动作轨迹之间仍然容易发生混淆。

后续可以通过引入更大规模的公开手语数据集、继续采集自有样本、增加不同用户和不同拍摄环境下的数据，
提升模型对个体差异、光照变化、动作速度和摄像角度变化的适应能力。

其次，
可以进一步扩大识别词表并优化词级分类模型。

当前系统的识别能力主要集中在固定小词表范围内，
适合用于毕业设计系统展示和短句识别验证，
但还不能覆盖真实交流场景中的开放词汇。

后续可以在现有特征工程基础上继续优化时序模型结构，
尝试引入更强的序列建模网络、注意力机制或预训练视觉特征，
提高模型对复杂动作和相似类别的区分能力。

同时，
也可以继续优化 Top-K 候选保留、词段合并、置信度过滤和 NMS 抑制策略，
减少插入误报、漏词和冗余输出问题。

再次，
可以继续探索更完整的连续手语识别方案。

本文在实验过程中尝试了基于 Transformer、Conformer 和 CTC 的连续识别方向，
但相关方案对数据规模、特征表示、训练策略和解码方法要求较高，
当前实验结果尚未达到稳定接入系统的程度。

后续如果能够获得更充分的连续手语标注数据，
并进一步掌握连续识别模型训练方法，
可以继续推进端到端连续手语识别研究，
使系统从“词级分类 + 滑动窗口后处理”逐步过渡到更接近真实连续手语识别任务的序列建模方案。

第四，
可以继续增强语义修正与自然语言表达能力。

当前 DeepSeek 语义修正模块采用受约束候选选择策略，
能够避免大语言模型自由生成缺乏视觉证据支持的内容，
但也限制了其对漏词和严重识别错误的修复能力。

后续可以进一步优化候选序列生成策略，
在保留视觉证据约束的前提下，
更充分地利用词段置信度、Top-K 候选、时间边界和语言合理性信息。

同时，
也可以探索更细粒度的提示词设计、候选重排序规则和错误解释机制，
使系统在输出自然中文表达的同时，
能够向用户说明结果的不确定性来源。

第五，
可以继续完善数字人动作资源和训练交互体验。

当前系统已经能够通过 SiGML 资源展示数字人动作，
但数字人资源覆盖范围和动作表现力仍然有限。

后续可以继续扩充手语动作资源库，
优化数字人动作播放体验，
并将数字人示范、用户实时识别结果、训练反馈和错误提示更紧密地结合起来，
形成更完整的交互式训练流程。

最后，
系统仍可在工程部署和真实用户评估方面继续完善。

当前系统主要在本地开发环境和毕业设计演示环境下运行，
后续可以进一步完善服务器部署、日志监控、异常告警、模型版本回滚、接口限流和数据备份机制。

同时，
也可以邀请听障用户、手语学习者或相关教师参与试用，
从真实用户角度评估系统在可理解性、易用性、训练帮助和交流辅助方面的实际效果。

通过持续优化数据、模型、交互和工程部署，
HearBridge 系统有望进一步发展为更加稳定、易用和具有实际应用价值的听障交流与训练辅助平台。